\documentclass[11pt]{article}

\usepackage[margin=2.2cm]{geometry}
\usepackage[utf8]{inputenc}
\usepackage[spanish]{babel}
\usepackage{amsmath, amssymb}
\usepackage{booktabs}
\usepackage{array}
\usepackage{xcolor}
\usepackage{tikz}
\usepackage{hyperref}

\definecolor{verde}{HTML}{1f7a1f}
\definecolor{rojo}{HTML}{b22222}
\definecolor{azul}{HTML}{0b3d91}

\newcommand{\W}{\mathbf{W}}
\newcommand{\Smat}{\mathbf{S}}
\newcommand{\Hen}{\mathcal{H}}

\title{Simetría de Hopfield: por qué los complementos son estados estables}
\author{SIA-TP4 — Aprendizaje No Supervisado}
\date{}

\begin{document}
\maketitle

\section*{Función de energía}

La energía de una configuración $\Smat \in \{-1, +1\}^N$ en una red de Hopfield con pesos $\W \in \mathbb{R}^{N \times N}$ se define como:

\[
\Hen(\Smat) \;=\; -\frac{1}{2}\sum_{i,\,j} W_{ij}\, S_i\, S_j
\]

Donde la doble suma recorre todos los pares de neuronas. Como $\W$ es simétrica ($W_{ij} = W_{ji}$) y la diagonal está en cero ($W_{ii} = 0$), cada par $(i,j)$ con $i\neq j$ aparece dos veces y contribuye con $W_{ij}\, S_i\, S_j$ en ambos sentidos.

\section*{Invariancia ante $\Smat \to -\Smat$}

Reemplazo $\Smat$ por su complemento $-\Smat$ (todos los píxeles invertidos):

\begin{align*}
\Hen(-\Smat)
  &= -\frac{1}{2}\sum_{i,j} W_{ij}\,(-S_i)(-S_j) \\
  &= -\frac{1}{2}\sum_{i,j} W_{ij}\,S_i\,S_j \\
  &= \Hen(\Smat).
\end{align*}

La doble negación se cancela: $(-1)\cdot(-1) = +1$. Por lo tanto $\Hen$ es \textbf{invariante} ante el cambio global de signo del estado.

\paragraph{Consecuencia.} Cualquier mínimo local $\Smat^{*}$ tiene un mínimo gemelo $-\Smat^{*}$ con exactamente la misma energía. El paisaje de energía es simétrico respecto del origen, como un par de pozos espejados.

\section*{Misma simetría en la regla de actualización}

La actualización síncrona es $\Smat_{t+1} = \operatorname{sign}(\W\,\Smat_t)$. Si $\xi$ es un estable, entonces $\operatorname{sign}(\W\,\xi) = \xi$, y:

\[
\operatorname{sign}\bigl(\W \cdot (-\xi)\bigr) \;=\; \operatorname{sign}(-\W\,\xi) \;=\; -\operatorname{sign}(\W\,\xi) \;=\; -\xi.
\]

Por cada patrón almacenado $\xi^{\mu}$, su complemento $-\xi^{\mu}$ también es estable. Con $p$ patrones almacenados, la red tiene $2p$ estados estables garantizados por esta simetría (más los espurios “mezcla” que aparezcan).

\section*{El término $W_{ij}\, S_i\, S_j$ en detalle}

El argumento de la simetría se entiende mejor mirando un único término de la suma. El producto $S_i\,S_j$ solo puede tomar dos valores, según si las neuronas $i$ y $j$ están en el mismo estado o no:

\begin{center}
\renewcommand{\arraystretch}{1.3}
\begin{tabular}{>{\centering\arraybackslash}m{2.2cm} >{\centering\arraybackslash}m{2.2cm} >{\centering\arraybackslash}m{2.5cm} >{\centering\arraybackslash}m{4cm}}
\toprule
$S_i$ & $S_j$ & $S_i\,S_j$ & Término $W_{ij}\,S_i\,S_j$ \\
\midrule
$+1$ & $+1$ & $+1$ & $+W_{ij}$ \\
$-1$ & $-1$ & $+1$ & $+W_{ij}$ \\
\midrule
$+1$ & $-1$ & $-1$ & $-W_{ij}$ \\
$-1$ & $+1$ & $-1$ & $-W_{ij}$ \\
\bottomrule
\end{tabular}
\end{center}

\paragraph{Lectura.} Hay dos casos:

\begin{itemize}
\item Si $S_i$ y $S_j$ están \textcolor{verde}{\textbf{alineados}} (ambos $+1$ o ambos $-1$), el término aporta $+W_{ij}$.
\item Si están \textcolor{rojo}{\textbf{anti-alineados}} (uno $+1$, otro $-1$), aporta $-W_{ij}$.
\end{itemize}

El término solo depende de la \emph{compatibilidad relativa} entre $S_i$ y $S_j$, no de sus signos absolutos. Si das vuelta los dos factores a la vez ($S_i \to -S_i$ y $S_j \to -S_j$), $S_i\,S_j$ no cambia. Por eso, cuando complementás \emph{toda} la red, cada término $W_{ij}\,S_i\,S_j$ queda exactamente igual, y la suma global —que es $-2\,\Hen$— también.

\section*{Ejemplo concreto ($N = 4$)}

Un único patrón almacenado:

\[
\Smat^{*} \;=\; \begin{pmatrix} +1 \\ +1 \\ -1 \\ +1 \end{pmatrix},
\qquad
-\Smat^{*} \;=\; \begin{pmatrix} -1 \\ -1 \\ +1 \\ -1 \end{pmatrix}.
\]

Pesos construidos con la regla de Hebb $\W = \frac{1}{N}\,\Smat^{*}\,\Smat^{*\,\top}$ y diagonal puesta en cero:

\[
\W \;=\; \frac{1}{4}\begin{pmatrix}
0 & +1 & -1 & +1 \\
+1 & 0 & -1 & +1 \\
-1 & -1 & 0 & -1 \\
+1 & +1 & -1 & 0
\end{pmatrix}
\;=\;
\begin{pmatrix}
0 & +0{,}25 & -0{,}25 & +0{,}25 \\
+0{,}25 & 0 & -0{,}25 & +0{,}25 \\
-0{,}25 & -0{,}25 & 0 & -0{,}25 \\
+0{,}25 & +0{,}25 & -0{,}25 & 0
\end{pmatrix}.
\]

Calculando energía y aplicando la regla:

\begin{center}
\renewcommand{\arraystretch}{1.25}
\begin{tabular}{l c c}
\toprule
Estado & $\Hen(\cdot)$ & $\operatorname{sign}(\W\cdot\,\cdot)$ \\
\midrule
$\Smat^{*} = (+1,+1,-1,+1)$ & $\mathbf{-1{,}50}$ & $\Smat^{*}$ \quad\textcolor{verde}{(estable)} \\
$-\Smat^{*} = (-1,-1,+1,-1)$ & $\mathbf{-1{,}50}$ & $-\Smat^{*}$ \quad\textcolor{verde}{(estable)} \\
$(+1,-1,+1,-1)$ \;(arbitrario) & $0{,}00$ & no es estable \\
\bottomrule
\end{tabular}
\end{center}

Los dos estados (patrón y complemento) tienen la misma energía $-1{,}5$, y ambos son estados estables. \textbf{La única decisión humana fue almacenar $\Smat^{*}$}; el complemento aparece solo, como consecuencia de la simetría de $\W$.

\section*{Visualización del paisaje}

\begin{center}
\begin{tikzpicture}[scale=1.2]
  % Doble pozo simétrico: f(x) = 0.05*(x^2 - 5.76)^2 - 1.5
  % En x=±2.4 → f = -1.5 (fondo de los pozos)
  % En x=0   → f = 0.05*33.18 - 1.5 ≈ 0.16 (cima en el medio)
  \draw[thick, azul, smooth, samples=300, domain=-3.8:3.8]
    plot (\x, {0.05*(\x*\x - 5.76)*(\x*\x - 5.76) - 1.5});
  % Eje x (espacio de estados)
  \draw[->, thick] (-4.2, -1.8) -- (4.2, -1.8) node[right] {$\Smat$};
  % Eje y (energía)
  \draw[->, thick] (-4, -1.85) -- (-4, 1.5) node[above] {$\Hen(\Smat)$};
  % Marcador del cero del eje x
  \draw (0, -1.75) -- (0, -1.85) node[below] {$\mathbf{0}$};
  % Marcadores de los dos pozos (fondo de cuencas) en y = -1.5
  \filldraw[verde] (-2.4, -1.5) circle (3pt);
  \filldraw[verde] ( 2.4, -1.5) circle (3pt);
  \node[verde, below] at (-2.4, -1.5) {$-\Smat^{*}$};
  \node[verde, below] at ( 2.4, -1.5) {$\Smat^{*}$};
  % Línea punteada uniendo los dos pozos a la misma altura
  \draw[dashed, gray] (-2.4, -1.5) -- (2.4, -1.5);
  \node[gray, above] at (0, -1.55) {misma $\Hen$};
  % Punto del estado arbitrario en una loma cerca del centro
  % En x=0.7: f = 0.05*(0.49-5.76)^2 - 1.5 = 0.05*27.77 - 1.5 ≈ -0.11
  \filldraw[rojo] (0.7, -0.11) circle (3pt);
  \node[rojo, right] at (0.85, -0.05) {estado arbitrario};
  % Flechas indicando la cuenca de atracción
  \draw[->, thick, gray!70] (-1.2, -0.5) -- (-2.0, -1.3);
  \draw[->, thick, gray!70] ( 1.2, -0.5) -- ( 2.0, -1.3);
\end{tikzpicture}
\end{center}

Las dos cuencas son espejos una de la otra. Cualquier input con $p_{\mathrm{flip}} < 0{,}5$ cae en la cuenca de $\Smat^{*}$; con $p_{\mathrm{flip}} > 0{,}5$ cae en la cuenca de $-\Smat^{*}$. El input está más cerca, en Hamming, del complemento que del original.

\section*{Implicancia para los experimentos del TP}

En los CSV el complemento queda etiquetado como \texttt{outcome=Espurio}, porque la definición operativa de “espurio” es \emph{“no es ninguno de los patrones almacenados”}. Esto es correcto. Pero conviene tener presente:

\begin{itemize}
  \item El complemento no es un espurio “accidental”: aparece \emph{garantizado} por la simetría de $\Hen$.
  \item Para la consigna 2.1.b (identificar un estado espúreo) preferimos destacar un caso de \emph{mezcla}, con $\operatorname{hamming}$ intermedio respecto de todos los almacenados, en vez de un complemento puro (\texttt{hamming\_final = 25} contra el original).
  \item Si en alguna slide aparece un complemento, alcanza con mencionar la simetría de $\Hen$ sin necesidad de profundizar.
\end{itemize}

\end{document}
